% ----------------------------------------------------------
% 1. Glossário - Definições e Siglas
% ----------------------------------------------------------
\chapter{Glossário -- Definições e Siglas}
\label{cap:glossario}

\section*{Siglas e termos da organização e do negócio}

\begin{description}
    \item[\textbf{EJ}] Empresa Júnior --- associação civil sem fins lucrativos gerida por discentes de graduação, prestando serviços e projetos para aprimoramento profissional acadêmico.
    \item[\textbf{Commit Júnior}] Nome da Empresa Júnior dos cursos de Engenharia do CEFET-MG, atuando como organização cliente e contexto de aplicação do sistema.
    \item[\textbf{Squad}] Subgrupo ou time ágil de trabalho alocado dentro de uma diretoria ou gerência técnica da empresa.
    \item[\textbf{RG (Reunião Geral)}] Reunião plenária periódica e institucional obrigatória, convocada para todos os integrantes ativos da empresa júnior.
    \item[\textbf{1-on-1}] Reunião individual de alinhamento e acompanhamento realizada entre dois membros específicos, desvinculada de eventos plenários.
    \item[\textbf{Alumnus}] Status conferido a ex-membros desligados da empresa júnior, preservando seus registros históricos de atuação sem exclusão física do banco de dados.
    \item[\textbf{Agravante}] Fator multiplicador incidente sobre a pontuação de penalidades quando constatada reincidência infracional antes do prazo regular de expiração.
    \item[\textbf{SIGAA}] Sistema Integrado de Gestão de Atividades Acadêmicas --- plataforma oficial do CEFET-MG em que as horas de extensão validadas são registradas pela coordenação de curso.
    \item[\textbf{CPF}] Cadastro de Pessoa Física --- registro documental civil empregado como identificador primário para assegurar a unicidade cadastral no sistema.
    \item[\textbf{SMTP}] \textit{Simple Mail Transfer Protocol} --- protocolo de comunicação em rede utilizado pelo servidor institucional para o despacho automatizado de notificações e relatórios por correio eletrônico.
\end{description}

\section*{Siglas de engenharia de software e modelagem}

\begin{description}
    \item[\textbf{DCU}] Diagrama de Casos de Uso --- representação gráfica estruturada segundo o padrão UML para mapear atores, casos de uso e os limites de escopo funcional do sistema.
    \item[\textbf{UML}] \textit{Unified Modeling Language} --- notação gráfica padronizada para especificação, modelagem e documentação de artefatos de software orientados a objetos.
    \item[\textbf{UC}] Caso de Uso (\textit{Use Case}) --- identificador padronizado atribuído a cada funcionalidade ou serviço executado pela aplicação (ex.: UC-01, UC-02).
    \item[\textbf{RF}] Requisito Funcional --- declaração dos comportamentos, serviços e processamentos de dados que o software deve prover aos usuários.
    \item[\textbf{RNF}] Requisito Não-Funcional --- especificação de critérios de qualidade, restrições e atributos operacionais da solução (tais como segurança, tempo de resposta e portabilidade).
    \item[\textbf{RN}] Regra de Negócio --- preceito normativo, condição lógica ou restrição operacional inerente ao funcionamento da instituição que delimita o sistema.
\end{description}

\section*{Termos técnicos de segurança e arquitetura}

\begin{description}
    \item[\textbf{RBAC}] \textit{Role-Based Access Control} (Controle de Acesso Baseado em Papéis) --- paradigma de segurança em que os privilégios e permissões de acesso às funcionalidades são determinados estritamente pelo perfil funcional exercido pelo usuário.
    \item[\textbf{Hash criptográfico / BCrypt}] Algoritmo matemático unidirecional de derivação com salgamento seguro, utilizado para converter credenciais de senha em cadeias irreversíveis, prevenindo exposição mesmo em vazamentos de dados.
\end{description}

\section*{Terminologia da notação de Casos de Uso (UML)}

\begin{description}
    \item[\textbf{Ator}] Entidade externa ao ambiente computacional (como pessoa, setor ou subsistema de hardware) que interage ativamente com o sistema.
    \item[\textbf{Ator primário}] Papel que dispara e conduz ativamente a interação com o caso de uso visando atender a uma demanda própria.
    \item[\textbf{Ator secundário}] Entidade externa periférica acionada pelo sistema durante a execução do fluxo para auxílio ou registro, sem que tenha iniciado a requisição original.
    \item[\textbf{Caso de uso}] Unidade lógica de comportamento observável que um ator executa para alcançar um resultado mensurável de valor, abstraindo detalhes internos de código.
    \item[\textbf{Generalização (herança)}] Relação estrutural hierárquica na qual um elemento especializado reaproveita e estende atributos e comportamentos do elemento genérico antecedente.
    \item[\textbf{Inclusão (\guillemotleft include\guillemotright)}] Relacionamento de dependência obrigatória no qual a realização de um caso de uso incorpora incondicionalmente a rotina descrita em outro caso de uso acoplado.
    \item[\textbf{Extensão (\guillemotleft extend\guillemotright)}] Relacionamento condicional e opcional que viabiliza a execução complementar de fluxos específicos a partir de pontos de extensão demarcados no caso de uso base.
\end{description}